% エントリ 19(2026-11-14): 独立な査読(三度目)と訂正。約束の破れ(段差の時刻の折れ)、単体版の移植性、印 3 の範囲
\section{[E19, 2026-11-14] 独立な査読と訂正: 約束の破れ(段差の時刻の折れ)、単体版の移植性、印 3 の範囲}\label{sec:e19}

E18 の版の χ$_\mathrm{SR}$ と、ノートの式・計算コード・共同研究者に渡す単体版の全体を、独立に査読した(ユーザー指示「査読して」、2026-11-13)。見つかった問題は直した(「直して E19 に記録」、2026-11-14)。

結論を先に書く。式、保存則、率の規約、基準の運動論の実装、表の引き方に誤りはなく、窓($f\sim10^{12}\GeV$)の数値と論文の結論は変わらない。しかし E18 の約束(印 0・1・2 の値は基準の運動論と 0.03 以内)は、\textbf{段差の時刻が $\alpha$ の節のあいだで折れる所}で破れていた。新しい照合では、E18 の版は印 0・1・2 のまま最大 0.86 ずれる。原因は基準の運動論の側にあり、(a) 超放射の率が連分数法の表と解析式で切り替わる $\alpha$ での不連続と、(b) 322 の率の表の節の誤差(倍精度の限界)である。折れを表そのものから測って印 5 にする規則を加え、独立な 4{,}700 点 × 2382 年齢で、越えるセルは 0 になった。このほかに、単体版が x86 の計算機で止まりうること、自由場で 544 の超放射まで進んだ点が印 2 に戻っていたこと、印 3 が広すぎたこと、単体版が pickle できないことなどを直した。

計算は \texttt{calc/e19\_review.py}(所見の確かめと、修正の前後の比較。出力 \texttt{calc/out/e19.json})と \texttt{calc/e19\_validate.py}(照合。出力 \texttt{calc/out/e19\_validate\_*.json}、\texttt{e19\_dense\_*.json})。修正の前の版は git の 7846eae(E18)から読み込む。計算機は Apple M3(8 コア)で、E18 の保存してある出力を完全に再現する。

\subsection{確かめたこと(誤りはなかった)}

\begin{itemize}
\item \textbf{式}: ノートの主要な式を手で導き直した。E2 の率の規約(占有数の率は $2\,\mathrm{Im}\,\omega$)と $e$-fold の時間、E3 の準平衡の占有数、E9 の Bardeen の因子(1.63、1.074)と質量の補正 $1-2\chi\alpha/3$、E10 の $H$、E12 の固定点の二次式 $3u^2+(2\gamma_{322}-4\gamma_\infty\gamma_{211}/\gamma_\mathrm{BH})u-\gamma_{322}^2=0$、E14 の三準位の式と角運動量の収支。どれもコードと合う。
\item \textbf{基準の運動論}: E18 の式だけから別に書いた三準位の ODE(変数は $\ln n_i$ と、初期値で割った BH の質量・角運動量。時間は $1/\mu$ の単位)は、\texttt{fp3.evolve3m} と 12 例 × 9 年齢で $2\times10^{-8}$ 以内で一致する(\texttt{e19\_review.py ode})。表の値は同じ点でこの ODE と 0.004 以内。
\item \textbf{再現性}: 修正の前の単体版は、保存してある E18 の照合の出力(seed 1801・1802、20 万セル)を、別の計算機・別の numpy(2.3)で、値も印も完全に再現する。
\item \textbf{窓の数値}(E12・E13 の値を表の版の χ$_\mathrm{SR}$ で計算し直した。残っていた作業): 表~\ref{tab:e19window}。$f$ の上限は E12 の半解析の値より 2.5--3.0\% 大きく、回転を失う 1 世代目の最小質量は 0.3$\Msun$ 以内で同じ。上限のところの印は 1 で、$f_\mathrm{max}\sqrt{T_2}$ は $T_2=3\times10^7$--$10^9\yr$ で $10^{-4}$ の精度で一定($f^2T$ の縮退)。E13 の結論は変わらない。
\end{itemize}

\begin{table}[H]
\centering\small
\begin{tabular}{cccccc}
\toprule
$T_2$ [yr] & $\chi_0$ & $f_\mathrm{max}$ [GeV](表の版) & E12(半解析) & 比 & 1 世代目の最小質量 [$M_\odot$]($\delta=0.02$、0.05)\\
\midrule
$5\times10^7$ & 0.78 & $1.91\times10^{12}$ & $1.85\times10^{12}$ & 1.030 & 13.3、11.9\\
$5\times10^7$ & 0.72 & $1.15\times10^{12}$ & $1.12\times10^{12}$ & 1.027 & 16.0、14.5\\
$1\times10^{8}$ & 0.78 & $1.35\times10^{12}$ & $1.31\times10^{12}$ & 1.028 & 15.1、13.6\\
$1\times10^{8}$ & 0.72 & $8.13\times10^{11}$ & $7.93\times10^{11}$ & 1.025 & 18.2、16.4\\
$3\times10^{8}$ & 0.78 & $7.78\times10^{11}$ & $7.55\times10^{11}$ & 1.030 & 18.5、16.7\\
$3\times10^{8}$ & 0.72 & $4.69\times10^{11}$ & $4.58\times10^{11}$ & 1.025 & 22.2、20.1\\
\bottomrule
\end{tabular}
\caption{GW241011 の主星($19.6\Msun$、$\mu=0.71\times10^{-12}\eV$)が年齢 $T_2$ のあと $\chi\ge0.69$ を保つ $f$ の上限と、その $f$ で $6\times10^9\yr$ のうちに 211 の臨界 $+\delta$ まで回転を失う 1 世代目($\chi_0=0.7$)の最小質量。}
\label{tab:e19window}
\end{table}

\subsection{所見 1: 段差の時刻が $\alpha$ の節のあいだで折れる所で、約束が破れていた}

\paragraph{原因 (a): 率の切り替えでの不連続}
基準の運動論は、超放射の率を連分数法の表(211 は $\alpha\ge0.03$、322 は $0.1\le\alpha\le0.52$)から引き、範囲の外は解析式(\arxiv{1908.10370} の式 2.32)を使う。二つは切り替えの $\alpha$ で一致しない(図~\ref{fig:e19}(a))。表の率は解析式の率に対して、322 の $\alpha=0.1$ で $+14\%$($\chi=0.25$)から $-31\%$($\chi=0.99$)、211 の $\alpha=0.03$ で $+2\%$ から $-12\%$ ずれる。$\alpha=0.1$ のすぐ上の表の値は多倍長の連分数法と 0.5\% 以内で合うので、ずれているのは解析式のほうである(所見 8)。

したがって基準の運動論は、$\alpha$ の関数として $\alpha=0.03$ と 0.1 で不連続である。値の表は $\alpha$ の節(0.02816 と 0.03081、0.09914 と 0.10846)のあいだを連続に補間するので、このセルの中では段差の時刻を誤る。BH の質量は超放射で減り、$\alpha$ は時間とともに下がるので、$\alpha_0$ が 0.1 より少し上の BH は進化の途中で切り替えをまたぐ。211 から育つ BH では、211 が飽和するまでに質量が最大 6\% 減るので、322 の段階の率が切り替わる $\alpha_0$ は初期スピンによって 0.100--0.106 のあいだを動く。

\paragraph{原因 (b): 322 の率の表の節の誤差}
E4 の 322 の表は倍精度の連分数法で作った。$\alpha\approx0.1$ の低いスピンでは $\mathrm{Im}\,\omega/\mathrm{Re}\,\omega\sim10^{-17}$ で、倍精度の丸め($10^{-16}$)より小さい。表の節の上の値を多倍長(50 桁)と比べると(153 節、\texttt{e19\_review.py noise322})、$\alpha\le0.125$ の低いスピンで節ごとにずれている。
\begin{center}\small
\begin{tabular}{cccc}
\toprule
$\alpha$ & $\chi$ & 表 / 多倍長 & $\mathrm{Im}\,\omega/\mathrm{Re}\,\omega$\\
\midrule
0.1018 & 0.21 & 1.079 & $4.2\times10^{-19}$\\
0.1195 & 0.27 & 0.973 & $1.1\times10^{-17}$\\
0.1195 & 0.30 & 1.092 & $2.1\times10^{-17}$\\
0.1243 & 0.27 & 1.011 & $1.4\times10^{-17}$\\
\bottomrule
\end{tabular}
\end{center}
ほかの節は 0.5\% 以内、$\alpha\ge0.146$ の節は 0.04\% 以内で合う。E4 で 322 を多倍長と比べたのは $(\alpha,\chi)=(0.15,0.9)$ の 1 点だけだった。433 の表が同じ理由で使えないことは E17 で見つけていたが、322 の低い $\alpha$ の端は確かめていなかった。率そのものへの影響は 1 割以下で物理には効かないが、表を補間した率が $\alpha$ について波打つ: 値の表のセルの真ん中での $\log_{10}(1/\gamma_{322})$ は、セルの両端を結ぶ直線から、$\alpha=0.108$--0.119 のセルで +0.017 dex(4\%)、0.119--0.130 で +0.007 dex 外れる(臨界のあいだの位置 $s=0.32$。解析式を使う $\alpha<0.099$ と、$\alpha>0.130$ では $3\times10^{-4}$ dex 以下)。

\paragraph{折れの大きさ}
値の表は、進み具合の準位ごとに、その準位に届く時刻 $t=\log_{10}\tau$ を持つ(E18)。$x=\ln\alpha$ の節 $x_i$ での値を $t_i$ と書く。セル $[x_i,x_{i+1}]$ について、下の 2 節を結ぶ直線を $x_{i+1}$ まで延ばした値と $t_{i+1}$ の差 $d_1$、上の 2 節を結ぶ直線を $x_i$ まで延ばした値と $t_i$ の差 $d_2$ を
\begin{align}
d_1&=\left|t_{i+1}-\Bigl[t_i+(t_i-t_{i-1})\frac{x_{i+1}-x_i}{x_i-x_{i-1}}\Bigr]\right|,&
d_2&=\left|t_i-\Bigl[t_{i+1}+(t_{i+1}-t_{i+2})\frac{x_i-x_{i+1}}{x_{i+1}-x_{i+2}}\Bigr]\right|
\label{eq:e19d}
\end{align}
と定義する。$t(x)$ がなめらかなら、$d_1$ も $d_2$ も二階の差分で小さい。セルの中に折れ(不連続)があれば、どちらの直線も折れをまたいで延ばすことになるので、両方が大きい。折れが隣のセルにあれば、片方だけが大きい。

表から測った $\max(d_1,d_2)$ と、折れとみなす準位($\min(d_1,d_2)$ が約束の許す時刻のずれ 2\% $=0.0087$ dex を越える準位)の割合は、次のとおりである(\texttt{e19\_review.py kink}。$f$ は GeV)。
\begin{center}\footnotesize
\begin{tabular}{llcccc}
\toprule
$\alpha$ のセル & 育ち方、$f$ & \multicolumn{3}{c}{$\max(d_1,d_2)$ [dex]} & 折れの準位 [\%]\\
 & & 中央値 & 90\% 点 & 99\% 点 & \\
\midrule
322 の切り替え(0.099--0.108) & 322 から、$f\ge10^{14}$ & 0.0920 & 0.142 & 0.196 & 100\\
322 の切り替え(0.099--0.108) & 211 から、$f\le10^{14}$ & 0.0005 & 0.001 & 0.003 & 0\\
322 の切り替え(0.099--0.108) & 211 から、$f\ge10^{14}$ & 0.0055 & 0.076 & 0.160 & 46\\
211 の切り替え(0.0282--0.0308) & 211 から、$f\le10^{14}$ & 0.0162 & 0.028 & 0.064 & 72\\
211 の切り替え(0.0282--0.0308) & 211 から、$f\ge10^{14}$ & 0.0121 & 0.043 & 0.061 & 66\\
その上(0.108--0.119) & 322 から、$f\ge10^{14}$ & 0.0846 & 0.142 & 0.196 & 65\\
その上(0.119--0.130) & 322 から、$f\ge10^{14}$ & 0.0280 & 0.096 & 0.146 & 49\\
解析式の範囲(0.076--0.083) & 322 から、$f\ge10^{14}$ & 0.0016 & 0.003 & 0.004 & 0\\
解析式の範囲(0.076--0.083) & 211 から、$f\le10^{14}$ & 0.0005 & 0.001 & 0.002 & 0\\
解析式の範囲(0.076--0.083) & 211 から、$f\ge10^{14}$ & 0.0025 & 0.008 & 0.076 & 8\\
表の範囲(0.155--0.170) & 322 から、$f\ge10^{14}$ & 0.0027 & 0.005 & 0.005 & 0\\
表の範囲(0.155--0.170) & 211 から、$f\le10^{14}$ & 0.0003 & 0.003 & 0.009 & 1\\
表の範囲(0.155--0.170) & 211 から、$f\ge10^{14}$ & 0.0044 & 0.006 & 0.263 & 4\\
\bottomrule
\end{tabular}
\end{center}
322 の切り替えのセルでは、322 から育つ BH の段差の時刻が、すべての準位で 0.07--0.2 dex(17--60\%)折れている。その上の二つのセルも、65\%、49\% の準位が折れている(原因 (b))。解析式だけを使う $\alpha$ と、表が正確な $\alpha\ge0.155$ では、322 から育つ BH に折れはない。同じ切り替えのセルでも、211 から育つ BH の窓の側($f\le10^{14}\GeV$)は折れていない(中央値 0.0005 dex)。降下の速さが 211 の率で決まり、322 の率にはほとんどよらないからである。211 から育つ BH の $f\ge10^{14}\GeV$ の数\% の折れは、どのセルにもある速い段差(E18 の印 5 が扱う所)である。

\paragraph{帰結}
図~\ref{fig:e19}(b)・(c) は、322 から育つ BH($60\Msun$、$\chi_0=0.29$、$f=5\times10^{18}\GeV$)で、$\mu$ を変えて $\alpha$ を動かした例である(\texttt{e19\_review.py stepscan})。切り替えのセルの中では、表の補間は段差の時刻を最大 11\% 遅く見積もり、E18 の版は印 2 のまま最大 0.092 ずれる。その上の $\alpha=0.108$--0.130 のずれ(-4\% から +5\%)は原因 (b) によるもので、E18 の版の誤差は最大 0.048。ほかの $\alpha$ では時刻は 0.2\% 以内で合う。

322 から育つ BH($f=10^{15}$--$10^{20.5}\GeV$、$M$ 2--200$\Msun$)を $\alpha$ のセルごとにとり、細かい年齢で照合すると(\texttt{e19\_review.py grown322})、E18 の版が約束を越える点は次のとおりである。
\begin{center}\footnotesize\setlength{\tabcolsep}{3.5pt}
\begin{tabular}{lccccccc}
\toprule
$\alpha$ のセル & 0.091--0.099 & 0.099--0.108 & 0.108--0.119 & 0.119--0.130 & 0.130--0.142 & 0.155--0.170 & 0.204--0.223\\
\midrule
点の数 & 235 & 244 & 244 & 246 & 255 & 257 & 242\\
E18 の版: 越えた点 & 0 & 82 & 1 & 3 & 0 & 0 & 0\\
E18 の版: 誤差の最大 & 0.002 & 0.177 & 0.039 & 0.046 & 0.004 & 0.006 & 0.011\\
E19 の版: 越えた点 & 0 & 0 & 0 & 0 & 0 & 0 & 0\\
E19 の版: 誤差の最大 & 0.002 & 0.015 & 0.014 & 0.015 & 0.004 & 0.006 & 0.011\\
印 5 の割合 [\%](E18 → E19) & 3.1 → 5.4 & 2.8 → 7.5 & 3.8 → 7.4 & 5.4 → 8.2 & 4.7 → 7.2 & 4.4 → 6.8 & 2.6 → 3.7\\
\bottomrule
\end{tabular}
\end{center}
切り替えのセルでは、322 から育つ自由場の側の BH の 3 分の 1 が、どこかの年齢で約束を越えていた。

\begin{figure}[H]
\centering
% E19: 超放射の率の切り替えと、段差の時刻の折れ。データは calc/e19_review.py rates(figs/data/e19_jump.dat、e19_step.dat)
\begin{tikzpicture}
\begin{groupplot}[group style={group size=3 by 1, horizontal sep=1.45cm}, note, width=0.345\linewidth, height=0.33\linewidth,
  title style={font=\small, yshift=-3pt}]
\nextgroupplot[title={(a) 切り替えでの率の飛び}, xlabel={BH のスピン $\chi$}, ylabel={表の率 / 解析式の率}, xmin=0.2, xmax=1.0, ymin=0.65, ymax=1.2,
  legend style={font=\scriptsize, at={(0.03,0.05)}, anchor=south west, fill=white, fill opacity=0.85, text opacity=1, draw=none}, legend cell align=left]
  \addplot[ink2, line width=0.4pt, forget plot] coordinates {(0.2,1) (1.0,1)};
  \addplot[cA, line width=1.0pt] table[x=chi, y=r322] {figs/data/e19_jump.dat}; \addlegendentry{322、$\alpha=0.1$}
  \addplot[cB, line width=1.0pt] table[x=chi, y=r211] {figs/data/e19_jump.dat}; \addlegendentry{211、$\alpha=0.03$}
\nextgroupplot[title={(b) 段差の時刻: 表の補間 / 基準}, xlabel={$\alpha$}, ylabel={時刻の比}, xmin=0.0915, xmax=0.128, ymin=0.92, ymax=1.16,
  xtick={0.10,0.11,0.12}, xticklabel style={/pgf/number format/fixed, /pgf/number format/precision=2}]
  \fill[cM!18] (axis cs:0.0915,0.98) rectangle (axis cs:0.128,1.02);
  \foreach \a in {0.09914,0.10846,0.11866} {\edef\temp{\noexpand\draw[ink2, dashed, line width=0.4pt] (axis cs:\a,0.92) -- (axis cs:\a,1.16);}\temp}
  \draw[cB, line width=0.6pt] (axis cs:0.1,0.92) -- (axis cs:0.1,1.16);
  \addplot[cA, line width=0.9pt, mark=*, mark size=0.9pt] table[x=alpha, y=ratio] {figs/data/e19_step.dat};
\nextgroupplot[title={(c) 印 0・1・2 の誤差の最大}, xlabel={$\alpha$}, ylabel={$|\Delta\chi|$(時刻の 2\% を許す)}, xmin=0.0915, xmax=0.128, ymin=0, ymax=0.2,
  xtick={0.10,0.11,0.12}, xticklabel style={/pgf/number format/fixed, /pgf/number format/precision=2},
  yticklabel style={/pgf/number format/fixed, /pgf/number format/precision=2},
  legend style={font=\scriptsize, at={(0.97,0.95)}, anchor=north east, fill=white, fill opacity=0.85, text opacity=1, draw=none}, legend cell align=left]
  \addplot[ink2, line width=0.4pt, dashed, forget plot] coordinates {(0.0915,0.03) (0.128,0.03)};
  \addplot[cB, line width=0.9pt, mark=*, mark size=0.9pt] table[x=alpha, y=err_old] {figs/data/e19_step.dat}; \addlegendentry{E18 の版}
  \addplot[cC, line width=0.9pt, mark=*, mark size=0.9pt] table[x=alpha, y=err_new] {figs/data/e19_step.dat}; \addlegendentry{E19 の版}
\end{groupplot}
\end{tikzpicture}

\caption{(a) 率の切り替えの $\alpha$ での、連分数法の表の率と解析式の率の比。(b) 322 から育つ BH($60\Msun$、$\chi_0=0.29$、$f=5\times10^{18}\GeV$)のスピンが 0.27 を切る時刻の、表の補間と基準の運動論の比。破線は値の表の $\alpha$ の節、橙の線は切り替え $\alpha=0.1$、灰色の帯は約束が許す $\pm2\%$。(c) 同じ BH の、印 0・1・2 の年齢での誤差の最大。破線は約束の 0.03。E19 の版は、ずれる年齢に印 5 をつける。}
\label{fig:e19}
\end{figure}

\paragraph{E18 の照合が見逃した理由}
E18 の照合は年齢を $10^5$--$1.4\times10^{10}\yr$ の対数で 40 点とった(隣り合う年齢の比は 1.35)。段差の時刻が 3--13\% ずれたときに約束を越える年齢の幅は、その 3--13\% しかないので、40 点の年齢がそこに当たる確率は 1 点あたり 1--4 割である。しかも折れのあるセルは $\alpha$ の 46 セルのうち数個で、E18 の層($\mu$ が対数一様、または 0.6--1.0$\times10^{-12}\eV$)ではそこに入る点が少なかった。E18 の入力(seed 1802、12 万セル)を修正の前の版で引くと、越えたセルは実際に 0 である。E18 の「どの範囲にも偏りはない」は、標本が足りない結論だった。

\paragraph{修正: 折れを表から測って印 5 にする}
基準の運動論の不連続をなくすには、率の表を多倍長で作り直して(322 を $\alpha=0.03$--0.15、211 を $\alpha<0.03$)、値の表を作り直す必要がある(多倍長の計算と、値の表に約 5 時間)。これは残る作業とし、ここでは約束を守る印をつける。

セル $[x_i,x_{i+1}]$ の中で、本当の時刻 $t(x)$ が、折れの位置 $x_e$ より下では下の直線の延長 $t_A(x)=t_i+(t_i-t_{i-1})(x-x_i)/(x_i-x_{i-1})$、上では上の直線の延長 $t_B(x)$ の上にあるとする。表の補間は $t_\mathrm{int}(x)=(1-u)t_i+u\,t_{i+1}$、$u=(x-x_i)/(x_{i+1}-x_i)$。$t_\mathrm{int}$ と $t_A$ はどちらも $x$ の一次式で、$x_i$ で一致するから、差は $u$ に比例し、$x_{i+1}$ での差が $\pm d_1$ である。上の側も同じで、
\begin{equation}
|t_\mathrm{int}(x)-t(x)|=\begin{cases}u\,d_1&(x<x_e)\\(1-u)\,d_2&(x>x_e)\end{cases}\ \le\ \max(d_1,d_2).
\label{eq:e19bound}
\end{equation}
つまり補間の時刻の誤差は $D\equiv\max(d_1,d_2)$ 以下である。規則は次のとおり(\texttt{reftab.compute\_ks}、\texttt{lookup\_rt3})。
\begin{enumerate}
\item 表を読むときに、$\alpha$ のセル・$s$ と $f$ の節・$M$ の節・準位ごとに $d_1$、$d_2$ を測る。$\min(d_1,d_2)>\delta_t$ の準位を折れとみなして $D$ を持つ($\delta_t=0.02/\ln10=0.0087$ dex は約束が許す時刻のずれ。$\min$ をとるのは、折れが隣のセルにある場合を除くため)。両方の節が届いている準位だけ測り、外の節が届いていなくて測れない準位は上限とする。準位 4 つずつの最大を 0.002 dex 刻みで持つ。
\item 問いの点について、そのセルの $D$ の最大(角の $s$・$f$ の節、$M$ の両方の節、問いの時刻に角の節がいる準位とその次の準位)を $D_q$ とする。
\item $D_q>\delta_t$ なら、時刻を $10^{\pm\min(D_q,0.3)}$ 倍にずらして表を引き直す。値が $\theta=0.015$ より変われば印 5(印 1・2 の点。値は補間のまま返す)。
\end{enumerate}
印がつかない点では、時刻が $D_q$ ずれても値は $0.015$ しか変わらないので、誤差は $0.015$ とふつうの補間の誤差(0.01 以下)の和で、0.03 に収まる。閾値は約束そのものから決めた(時刻の 2\%、値の 0.03 の半分)。標本に合わせた調整はしていない。E18 の印 5(角の節の値の幅 $>0.06$ かつ曲がり $>0.05$ dex)はそのまま残す。E18 の曲がりの指標は両隣を結ぶ直線からのずれで、0.05 dex(12\%)より小さい折れを拾わない。

図~\ref{fig:e19} の例では、E19 の版の印 0・1・2 の誤差は最大 0.0008 になる。

\paragraph{段のあとの平らな所の値}
折れの印を入れたあとに残った誤差の大きい点(0.02 以上。約束の中だが、余裕が小さい)は、どれも段のあとの平らな所の値のずれだった。原因は表の引き方にある。表は、進み具合の準位ごとの時刻を節のあいだで混ぜるので、重みのある角の節のうち一つでも次の準位に届かないと、補間はそこで止まり、値は進みのいちばん遅い節の値になる。その節の重みが小さくても同じである。

例(\texttt{e19\_review.py plateau}): $\alpha=0.261$、$24.8\Msun$、$\chi_0=0.862$、$f=10^{17.498}\GeV$ の BH は、211 と 322 が飽和したあと、BH の質量が減って下がった 322 の臨界(0.461。進み具合 $b=$ 2.08)で止まる。$f$ の節 $10^{17.5}\GeV$ の側の角は $b=$ 2.06--2.09 まで届くが、$10^{17.25}\GeV$ の側の角は $b=$ 2.02 で止まる。こちらの重みは 0.6\% しかないのに、表は $b=$ 2.02 の値 0.481 を返す(誤差 0.020)。

本当の値は角の節の値のあいだにあるので、誤差は角の節の値の幅以下である。補間が止まった所で、角の節の値の幅が $\theta_p=0.02$ を越える点を印 5 にした。この例は印 5 になる。引き方そのものを直す(止まった所の値を、角の節の値の重みつき平均にする)のは、値が変わるので残る作業とする。

結果は「修正のあとの照合」の節にまとめる。

\subsection{所見 2: 単体版が x86 の計算機で止まりうる}

表を引く関数(\texttt{\_core2})は、重みが 0 の節でも戻りの表の添字を計算する。その節がまだ準位に届いていないと時刻が NaN になり、NaN を整数に直していた。この結果は CPU で違う。Apple Silicon では 0 になって結果に影響しないが(重みが 0 なので)、x86 では $-2^{63}$ になり、配列の添字として範囲の外なので \texttt{IndexError} で止まる。Mac Studio と MacBook(どちらも Apple Silicon)でしか走らせていなかったので、見えていなかった。

x86 の挙動(NaN を $-2^{63}$ にする)を模した版で確かめると(\texttt{e19\_review.py usability})、E18 の版は、$f=10^{10}$--$10^{19}\GeV$ と $10^{17}$--$10^{20}\GeV$ の 20 万個の標本(5--60$\Msun$、$\chi_0$ 0--0.99、$\mu$ 0.6--1.0$\times10^{-12}\eV$、年齢 $10^5$--$10^{10}\yr$)で \texttt{IndexError} になった。窓の $10^{11}$--$10^{13}\GeV$ だけの標本では起きない。NaN を 0 に置き換えてから整数に直すようにした。値は変わらず(E18 の照合の出力を完全に再現)、模した版でも NaN が整数に直される回数は 0 になった。実機の x86 では走らせていない。渡した先で値と印を照らす \texttt{share/chi\_sr/selftest.py} を付けた(いろいろな領域の 5000 個の入力と、ここでの出力。警告を例外にして呼ぶ)。

\subsection{所見 3: 印 3 が広すぎた($\alpha<0.01$ の固定点)。値の表の下の扱い}

E18 の版は、固定点の領域で $\alpha<0.01$ の点を一律に印 3 にしていた(印 4 の判定に使う半解析の表が $\alpha=0.01$--0.30 なので)。E18 の照合の入力(seed 1802)では、印 3 のセルの 47\% がこれである。$\mu=10^{-13}\eV$ では $\alpha<0.01$ は $13\Msun$ 以下にあたり、5--10$\Msun$ の 92\%、10--20$\Msun$ の 37\% が印 3 だった(1 世代目の自然スピン、$f=10^{11}$--$10^{13}\GeV$)。低い $\mu$ の端で、何も起きない軽い BH を「模型の外」として落とすことになる。

\begin{itemize}
\item \textbf{$0.008\le\alpha<0.01$、$f\ge10^8\GeV$}: 値の表が使えるので、表から引く。印 4 の判定の表 $J$ は $\alpha\le0.038$ で 0(三準位の外の準位は育たない)なので、$\alpha=0.01$ の行を使えば足りる。
\item \textbf{値の表の下}($\alpha<0.008$。固定点では $\alpha<0.01$ で $f<10^8\GeV$ の点も): 雲が 1 個の量子から固定点の占有数 $\varepsilon_\mathrm{eq}$ まで育つ時刻 $t_\mathrm{grow}=\ln(\varepsilon_\mathrm{eq}GM^2)/\gamma$ の 7 $e$-fold 前より若い点、つまり $\gamma(t_\mathrm{grow}-T)\ge7$ の点は、$\chi_0$ をそのまま返す(印 1)。それまでに雲が持つ角運動量は $e^{-7}\varepsilon_\mathrm{eq}$ 以下なので、失うスピンは $10^{-3}$ 以下である。そうでない点は印 3。自由場の側も同じにした(飽和の時刻 $\ln[(\chi_0-\chi_c)GM^2]/\gamma$ の 7 $e$-fold 前より若ければ印 2 で $\chi_0$、そうでなければ印 3)。E18 の版は、自由場の側では半解析の飽和の時刻を返していたが、これは約束の 2\% より粗いことがある(下の確かめ)。
\end{itemize}

確かめ(\texttt{e19\_review.py lowalpha}): $\alpha=0.005$--0.008 で、雲が年齢のうちに育ちうる軽い BH(1--3.5$\Msun$)の 500 点を、細かい年齢で照合した。基準の運動論が 0.01 より動くのは 18 点で、どれも $2\Msun$ 未満である。E19 の版は、印 1 の 694{,}722 セルと印 2 の 498{,}901 セルで誤差の最大が $6\times10^{-5}$(基準は動いていない)。基準が 0.01 より動いたセルは、すべて印 3 である(印 3 は全体の 0.2\%)。E18 の版は、58\% が印 3 で(そのうち基準が 0.01 より動いたのは 1 セル)、印 2 の 9 セル(7 点、どれも $2\Msun$ 未満)が約束を越えていた(最大 0.053)。$2\Msun$ 以上でも、雲が年齢のうちに育つ 2--3$\Msun$ の隅では同じことが起きうる。$\mu=10^{-13}\eV$ の 5--20$\Msun$ の印 3 は 0\% になった(所見 7 の計算)。

固定点で $\alpha>0.30$ の点(E18 の照合の入力で印 3 の 51\%)は、印 4 の判定をこの範囲に延ばす必要があり、印 3 のまま残す。

\subsection{所見 4: 544 の超放射まで進んだ自由場の点が、印 2 に戻っていた}

seed 1902 の照合で、折れの印を入れたあとも値が大きくずれる点が見つかった($\alpha=0.491$、$32.5\Msun$、$\chi_0=0.793$、$f=6\times10^{20}\GeV$)。照合の年齢の範囲($1.4\times10^{10}\yr$ まで)では E18 の印 5 がついているが、そのすぐ先($1.43\times10^{10}\yr$)で印 2 に戻り、誤差は 0.046 だった。段の時刻が少し早い点なら、範囲の中で同じことが起きる。調べると、折れとは別の問題だった(\texttt{e19\_review.py m4stage})。

この BH は 322 が飽和して 322 の臨界の近く(0.781)で止まり、三準位の運動論では $3.5\times10^{9}\yr$ ごろに 544($m=4$)が超放射で飽和して 0.440 まで下がる。E17 の 433 の印は「返すスピンから 433 が年齢のうちに飽和しうるか」を見る。322 の臨界で止まっているあいだは印 4 がつくが、544 の段階のあとは返すスピンが 433 の臨界(0.588)より下なので、印がつかずに印 2 に戻っていた。しかし止まっているスピンでの 433 の率は 544 の率の $2.7\times10^{4}$ 倍で、433 は $1.2\times10^{5}\yr$ で飽和する。544 の段階まで進んだ点は、その前に 433 が飽和しているはずで、三準位の値は使えない。

自由場(印 2)で返すスピンが 433 の臨界より下の点を、印 4 にした。この点の値のずれが大きいのは、段のあとの平らな所の値が、角の節のうち進みのいちばん遅い節で決まるためである(所見 1 の最後)。

\subsection{所見 5: 単体版の使い勝手}

\begin{itemize}
\item \textbf{pickle できなかった}: \texttt{ChiSR} が、開いたままの表のファイルを属性に持っていた(使っていない)。multiprocessing や、それを使うサンプラーに渡せない。ファイルは読み終えたら閉じ、pickle には表の場所だけを入れるようにした(129 バイト。受け取ったプロセスが同じ表を一度だけ読む)。
\item \textbf{メモリ}: 表を実数に直して持っていたので、1 個の \texttt{ChiSR} が 1.4 GB を使っていた(README になかった)。$\log_{10}\tau$ は uint16 の符号、戻りは float16 のまま持ち、引いた要素だけを実数に直すようにした。値の表の配列は 1.42 GB から 0.40 GB に、表を読むときのメモリの最大は 3.3 GB から 1.0 GB になった。値は変わらない(プロジェクト版は実数で持つので、両者の完全な一致がこの確かめになる)。
\item \textbf{警告}: 配列で呼ぶと \texttt{divide by zero}(始めの準位の $\log_{10}0$)と \texttt{invalid value encountered in cast}(所見 2)の警告が出ていた。なくした。書き出しの確かめ(\texttt{export\_chisr.py})は、単体版の呼び出しをすべて、警告を例外にして行う。
\item \textbf{速さ}: $10^6$ 個の標本で 4.7、10.9 秒から 5.9、14.2 秒(1 世代目らしい入力、$\chi_0$ が一様の順。Apple M3)。折れの判定のぶん 3 割ほど遅い。表を読む時間は 1.2 秒から 2.5 秒。
\item \textbf{基準の表がない計算機}: プロジェクト版(\texttt{chisr.build})は基準の表(700 MB、git に入れていない)がないと止まっていた。プロジェクト版はもともと丸めた値($\log_{10}\tau$ を uint16、戻りを float16)を使うので、単体版に書き出した表(12 MB、git にある)から同じ数の表を作れる。ないときはそうするようにした。書き出しの確かめ(単体版との完全な一致など)は、どの計算機でも走る。
\end{itemize}

\subsection{所見 6: 関数の質量は生まれたときの質量}

関数の $M$ は生まれたときの質量だが、README に書いていなかった。観測された(超放射で減ったあとの)質量を渡すと、$\alpha$ が小さくなる。基準の運動論で初期の質量から最後の質量とスピンを出し、最後の質量で関数を呼んで比べた(\texttt{e19\_review.py mass}。$M$ 8--60$\Msun$、$\mu$ 0.3--1.2$\times10^{-12}\eV$、$\alpha\le0.3$、年齢 $10^7$--$6\times10^9\yr$、どちらの呼び方でも印 0・1・2 のセル)。
\begin{center}\small
\begin{tabular}{lcccc}
\toprule
$f$ [GeV] & 質量の減少(中央値、最大) & ずれの 95\% 点 & ずれの最大 & 0.03 を越える割合\\
\midrule
$10^{11}$--$10^{13}$ & 0.00\%、2.0\% & 0.0039 & 0.008 & 0.0\%\\
$10^{14}$--$10^{16}$ & 0.48\%、2.6\% & 0.0064 & 0.012 & 0.0\%\\
$10^{17}$--$10^{20}$ & 1.69\%、11.1\% & 0.0212 & 0.038 & 1.3\%\\
\bottomrule
\end{tabular}
\end{center}
窓では無視できる。$f\gtrsim10^{17}\GeV$ では、質量が大きく減る点($\alpha=0.23$--0.25、質量の減少の中央値 11\%)で 0.03 を越える。README に定義と表を書いた。

\subsection{所見 7: 約束の外の標本の割合と、その扱い}

印 3・4・5 の標本の割合を、$\mu$・質量・スピンの分布ごとに数えた(\texttt{e19\_review.py coverage}。年齢は $10^7$--$10^{10}\yr$ の対数一様、1 セル 4 万個)。表~\ref{tab:e19cov} は $\mu=0.71\times10^{-12}\eV$ の場合である。
\begin{table}[H]
\centering\small
\begin{tabular}{llccccc}
\toprule
スピン & $f$ [GeV] & 5--10$\Msun$ & 10--20 & 20--40 & 40--80 & 80--150\\
\midrule
1 世代目 & $10^{11}$--$10^{13}$ & 0.0/0.0/0.0 & 0.0/0.0/0.1 & 0.0/3.2/0.1 & 6.9/0.8/0.0 & 76.8/0.0/0.0\\
速く回る群 & $10^{11}$--$10^{13}$ & 0.0/0.0/0.3 & 0.0/0.0/2.2 & 0.0/43.7/1.8 & 50.0/16.5/0.0 & 100.0/0.0/0.0\\
1 世代目 & $10^{16}$--$10^{19}$ & 0.0/0.0/0.1 & 0.0/0.0/5.0 & 0.0/1.8/3.0 & 4.9/3.8/0.1 & 76.4/0.4/0.0\\
速く回る群 & $10^{16}$--$10^{19}$ & 0.0/0.0/0.0 & 0.0/0.0/16.1 & 0.0/7.1/10.6 & 23.1/34.8/2.1 & 92.9/7.1/0.0\\
\bottomrule
\end{tabular}
\caption{約束の外の標本の割合 [\%](印 3 / 4 / 5 の順)、$\mu=0.71\times10^{-12}\eV$。1 世代目は Beta(0.69, 2.28)、速く回る群は $\chi_0$ が 0.6--0.9 の一様。}
\label{tab:e19cov}
\end{table}
\begin{itemize}
\item 1 世代目の自然スピンでは、$40\Msun$ までは約束の外は 5\% 以下である。80--150$\Msun$ は 77\% が印 3($\alpha>0.42$)。
\item 速く回る BH は、20--40$\Msun$ の 44\% が印 4、40--80$\Msun$ の 67\% が印 3・4 である。DESIGN.md D2 の 2 つ目の較正(45$\Msun$ 以上の速く回る群)は、ほとんどが約束の外にある。
\item $\mu$ が大きいと、重い BH から順に印 3 になる($\alpha>0.5$ か、固定点で $\alpha>0.30$)。1 世代目で印 3 の割合は、$\mu=1.5\times10^{-12}\eV$ で 20--40$\Msun$ の 9\%、40--80$\Msun$ の 86\%、$\mu=3\times10^{-12}\eV$ で 20--40$\Msun$ の 86\%。$(\mu,f)$ の地図がこの模型で決まるのは、およそ $\mu\lesssim10^{-12}\eV$ である。
\end{itemize}
印のついた標本を落として規格化し直すと、重い BH・速く回る BH ほど落ちるので、$(\mu,f)$ に依存した偏りが出る。README に、落とさずに割合を記録すること、「スピンは変わらない」と「322 の臨界まで落ちる」の両端で挟むことを書いた。

\subsection{所見 8: 322 の解析式の率は、臨界のあいだでは小さい(SESSION の記述の訂正)}

SESSION.md の残る作業は「322 の解析式は 15--25\% 大きめで、211 の臨界より下と 322 から育つ BH の降下を 25--35\% 速く見積もる」と書いていた。多倍長の連分数法と比べると(\texttt{e19\_review.py rates})、向きが逆である。
\begin{center}\small
\begin{tabular}{cccccc}
\toprule
 & \multicolumn{2}{c}{322 と 211 の臨界のあいだ} & $\chi=0.5$ & 0.7 & 0.9\\
$\alpha$ & 下から 1/4 & 3/4 &  &  & \\
\midrule
0.05 & 0.958 & 0.966 & 1.024 & 1.074 & 1.158\\
0.07 & 0.927 & 0.941 & 1.006 & 1.076 & 1.199\\
0.09 & 0.890 & 0.914 & 0.976 & 1.065 & 1.225\\
0.10 & 0.871 & 0.899 & 0.957 & 1.054 & 1.234\\
\bottomrule
\end{tabular}
\end{center}
(解析式 / 多倍長。)降下が起きる臨界のあいだのスピンでは、解析式は 3--13\% \textbf{小さい}。固定点の降下の速さは $\gamma_{322}^{3/2}$ に比例するので、$\alpha<0.1$ の 322 から育つ BH と 211 の臨界より下への降下を 5--20\% \textbf{遅く}見積もっている。解析式が大きいのは $\chi\gtrsim0.7$ の高いスピンで、「15--25\% 大きめ」はそこでの値だった。211 の解析式は、$\alpha=0.02$--0.03 で多倍長と -2\% から +14\% の範囲で合う(高いスピンほど大きい)。窓には効かない。

\subsection{所見 9: 422 を調べていなかった}

E15 は三準位の外の準位として 411・433・544 を調べ、422 は「411 のあとに育つ準位で、率は表にない」として扱わなかった。しかし 422 は 411 を経なくても、$211\times211\to422\times\mathrm{BH}$ で直接育つ。率は BGLS の Table V にあり、322 を育てる過程と同じ桁である。
\begin{equation}
\Gamma^{422\times\mathrm{BH}}_{211\times211}=1.5\times10^{-7}\alpha^{11}\tilde r_+,\qquad
\Gamma^{211\times\infty}_{322\times422}=1.6\times10^{-8}\alpha^{8}
\end{equation}
(どちらも $\times(\Mpl/f)^4\mu$。後者は \arxiv{2412.03655} の Table I。抜き書き \texttt{Refs/notes/theory\_rates.md} の 4.2 節・5 節の値で、\arxiv{2412.03655} の HTML 版の Table I から読み直した値と一致する。原文の PDF はこの計算機になく、BGLS の Table V は直接は見ていない。E15 が「表にない」とした $411\times411\to422\times\mathrm{BH}$ も \arxiv{2412.03655} の Table I にあり、$2.3\times10^{-7}\alpha^7$ である。)

\paragraph{(211, 322) の固定点}
422 の占有数について一次の項だけを残すと、1 粒子あたりの増加率は
\begin{equation}
\lambda_{422}=\gamma_{422}+\kappa'_\mathrm{BH}\alpha^{11}\tilde r_+\,\varepsilon_{211}^2-\kappa'_\infty\alpha^8\,\varepsilon_{211}\varepsilon_{322},
\qquad\kappa'_\mathrm{BH}=1.5\times10^{-7},\ \kappa'_\infty=1.6\times10^{-8}.
\end{equation}
固定点では(E12)、$u\equiv\gamma_\mathrm{BH}\varepsilon_{211}^2$、$\varepsilon_{322}=(u+\gamma_{322})/(2\gamma_\infty\varepsilon_{211})$ である($\gamma_\mathrm{BH}=\kappa_a\alpha^{11}\tilde r_+$、$\gamma_\infty=\kappa_b\alpha^8$ は 322 の過程の率。$(\Mpl/f)^4$ は占有数に吸収した)。代入すると
\begin{equation}
\lambda_{422}=\gamma_{422}-\Bigl(\frac{\kappa'_\infty}{2\kappa_b}-\frac{\kappa'_\mathrm{BH}}{\kappa_a}\Bigr)u-\frac{\kappa'_\infty}{2\kappa_b}\gamma_{322}
=\gamma_{422}-0.378\,u-0.727\,\gamma_{322}.
\label{eq:e19lam422}
\end{equation}
422 の超放射の率は、多倍長の連分数法で $\gamma_{422}=(0.60$--$0.62)\,\gamma_{322}$ である($\alpha=0.05$--0.3、20 点。水素原子の主要項では $C_{42}/C_{32}=0.60$)。$u>0$ なので
\begin{equation}
\lambda_{422}\le-0.11\,\gamma_{322}-0.378\,u<0
\end{equation}
で、\textbf{(211, 322) の固定点では 422 は育たない}(211 の臨界に近づくと $u\to\gamma_{322}/3$ で $\lambda_{422}\to-0.23\gamma_{322}$。$\alpha=0.03$--0.3 の格子でも正にならない)。ただし余裕は大きくない。BH が吸収する二つの過程の有限 $\alpha$ の補正が同じなら式~\eqref{eq:e19lam422} は変わらないが、$211\times211\to422\times\mathrm{BH}$ の補正が $211\times211\to322\times\mathrm{BH}$ の補正の 2.1 倍を越えると、$\lambda_{422}>0$ になる。\arxiv{2412.03655} は、422 が中間の $\alpha$ で $211\times211$ の散乱から育ちうると述べている。

\paragraph{(322, 544) の固定点}
211 が空なので、上の二つの過程は効かない。422 の超放射の率は正($0.60$--$0.62\,\gamma_{322}$)で、育つかどうかは 544 を含む過程(\arxiv{2412.03655} によれば $422\times422\to544\times\mathrm{BH}$ が 422 を抑える)の率で決まる。その率は抜き書きになく、\textbf{確かめていない}。E15 の「(322, 544) では $\alpha\le0.30$ で何も育たない」は、411 と 433 についての結果である。

\subsection{軽微な訂正}

\begin{itemize}
\item \texttt{calc/sr.py} の冒頭: \arxiv{1908.10370} の著者は Baumann, Chia, Stout, ter Haar(Porto ではない。抜き書きには正しく書いてあった)。
\item \texttt{calc/export\_chisr.py} の冒頭の「0.003 以内」は、コードとノート(E18)の 0.01 に合わせた。
\item \texttt{paper/outline.md} の残る作業の 3・4 に、E18 で済んだ内容が残っていたので直した。
\item CLAUDE.md の \texttt{\textasciitilde/miniforge3/bin/python3} は、MacBook では numpy を読めない。使える環境を書き足した。
\item \texttt{export\_chisr.py} は、表の中身が変わらないときはファイルを書き直さないようにした(zip の時刻だけが変わって git に差分が出るのを防ぐ)。
\end{itemize}

\subsection{修正のあとの照合}

\paragraph{新しい点、細かい年齢}
E18 と同じ約束・同じ基準の曲線の作り方で、新しい種を使い、E18 の 4 つの層(各 750 点)に、問題の見つかった所の層を加えた(各 600 点): $\alpha$ が 211 の切り替え(0.03)を挟むセルとその両側、322 の切り替え(0.1)を挟むセルとその両側、$\alpha=0.004$--0.012。加えた層は $M$ 2--200$\Msun$、$\chi_0$ 0--0.995、$f$ は半分を $10^6$--$10^{22}\GeV$、半分を $10^{16}$--$10^{20.5}\GeV$。基準の曲線は年齢の 2400 点の格子で持っているので、その格子の年齢のうち約束の範囲($10^5$--$1.4\times10^{10}\yr$)に入る 2382 点すべてで単体版を呼び、年齢の窓は格子の前後 4 点($|\ln(T'/T)|\le0.0199$)でとる。seed 1901 は折れの規則を作るあいだに見た。seed 1902 は折れの規則を確定したあとで作り、所見 4 を見つけた。seed 1903 は、コードをすべて確定したあとで作った独立な標本である。

\begin{table}[H]
\centering\footnotesize\setlength{\tabcolsep}{3.5pt}
\begin{tabular}{lcccccc}
\toprule
 & \multicolumn{2}{c}{seed 1901} & \multicolumn{2}{c}{seed 1902} & \multicolumn{2}{c}{seed 1903(独立)}\\
 & E18 の版 & E19 の版 & E18 の版 & E19 の版 & E18 の版 & E19 の版\\
\midrule
印 0・1・2 のセル & 10{,}106{,}686 & 10{,}508{,}767 & 10{,}078{,}502 & 10{,}481{,}010 & 10{,}174{,}029 & 10{,}567{,}480\\
印 0・1・2 を返す点 & 4{,}465 & 4{,}675 & 4{,}474 & 4{,}687 & 4{,}492 & 4{,}700\\
約束を越えたセル & 390 & 0 & 458 & 0 & 357 & 0\\
約束を越えた点 & 37 & 0 & 42 & 0 & 32 & 0\\
約束の誤差の最大 & 0.845 & 0.018 & 0.843 & 0.018 & 0.863 & 0.015\\
99\% 点 & 0.004 & 0.004 & 0.004 & 0.004 & 0.004 & 0.004\\
誤差 $>0.02$ のセル & 2{,}472 & 0 & 1{,}714 & 0 & 436 & 0\\
照合できなかった印 0・1・2 のセル & 357 & 357 & 0 & 0 & 0 & 0\\
印 3、4、5 の割合 [\%] & 6.3、2.7、2.6 & 1.9、2.7、3.5 & 6.7、2.5、2.6 & 2.2、2.6、3.6 & 6.1、2.3、2.6 & 1.7、2.3、3.5\\
印 4 の誤差(0.03 を越える割合、最大) & 0.0\%、0.020 & 0.0\%、0.034 & 0.2\%、0.139 & 0.6\%、0.139 & 0.0\%、0.016 & 0.0\%、0.016\\
印 5 の誤差(0.03 を越える割合、最大) & 2.0\%、0.80 & 1.6\%、0.85 & 3.1\%、0.55 & 2.1\%、0.84 & 2.3\%、0.72 & 1.8\%、0.86\\
\bottomrule
\end{tabular}
\caption{χ$_\mathrm{SR}$ の単体版と基準の運動論の照合(各 4800 点 × 2382 年齢)。誤差は年齢の 2\% のずれを許した値。}
\label{tab:e19val}
\end{table}

E18 の版は、どの種でも約束を越える。印 0・1・2 を返す点の 0.7--0.9\%(37、42、32 点、390、458、357 セル)が、どこかの年齢で 0.03 を越え、最大は 0.86。越える点は、率の切り替えのセルを狙った層だけでなく、E18 と同じ作り方の 4 つの層にもある(三つの種で 8 点。最大は一様の層 0.65、臨界の近くの層 0.19、1 世代目の層 0.17、GW241011 の層 0.037)。

E19 の版は、三つの種のどれでも約束を越えるセルがない(合わせて 31{,}557{,}257 セル、14{,}062 点)。誤差の最大は 0.015--0.018 で、約束の半分ほどである。印 3 は全体の 6.1--6.7\% から 1.7--2.2\% に減り、印 5 は 2.6\% から 3.5--3.6\% に増えた。層ごとの印 5 は、1 世代目の層で 0.7--0.9\% から 0.9--1.1\%、GW241011 の層で 7.0--7.7\% から 8.4--9.1\%(年齢 200 点の評価)。印 4・5 の値は予言として使えない(E19 の版で、印 5 の 1.6--2.1\% が 0.03 を越え、最大 0.86。印 4 は 0.0--0.6\%、最大 0.14)。

\paragraph{E18 の照合の入力}
E18 の照合の入力(年齢 40 点)を E19 の版で引き直した(\texttt{e19\_review.py e18})。値が変わらないセルは保存してある約束の誤差を、印 3 から値を返すようになったセルは $|\chi_\mathrm{SR}-\chi_\mathrm{ref}|$ を使う。
\begin{center}\small
\begin{tabular}{lcc}
\toprule
 & seed 1801 & seed 1802\\
\midrule
印 0・1・2 のセル(E18 の版 → E19 の版) & 69{,}004 → 70{,}917 & 104{,}941 → 108{,}128\\
約束の誤差の最大 & 0.021 → 0.016 & 0.023 → 0.023\\
約束を越えたセル & 0 → 0 & 0 → 0\\
印 3 → 1(値を返すようになった) & 2{,}400 & 3{,}715\\
印 1 → 5、2 → 5、3 → 5 & 109、378、0 & 186、342、5\\
印 3・4・5 の割合 [\%](E19 の版) & 3.6、4.8、3.0 & 3.5、3.5、2.9\\
\bottomrule
\end{tabular}
\end{center}

\paragraph{単体版とプロジェクト版}
無作為な 20 万個、E18 の照合の 12 万個の入力、E19 で印の付け方を変えた所($\alpha$ の切り替えのセルとその両側、$\alpha<0.012$)に置いた 10 万個、上の照合のすべてのセルで、値も印も完全に一致する。単体版の呼び出しは、すべて警告を例外にして行った。

\begin{result}
\begin{itemize}
\item 式・保存則・率の規約・基準の運動論の実装・表の引き方に誤りはない。窓の数値は表の版で $f$ が 3\% 大きくなるだけで、E13 の結論は変わらない。
\item E18 の約束は、段差の時刻が $\alpha$ の節のあいだで折れる所で破れていた(新しい照合の三つの種で、印 0・1・2 を返す点の 0.7--0.9\% がどこかの年齢で越え、最大 0.86)。原因は、率の切り替えでの基準の運動論の不連続と、322 の率の表の節の誤差(倍精度の限界、最大 9\%)。折れを表から測って印 5 にした。独立な標本で、E19 の版が約束を越えるセルは 0(独立な seed 1903 の 4{,}700 点、10{,}567{,}480 セル。誤差の最大 0.015)。
\item 同じ照合で見つかった二つの穴もふさいだ: 段のあとの平らな所で角の節の値が割れている点(印 5)と、自由場で 544 の超放射まで進んだ点(印 4)。
\item 単体版は x86 の計算機で止まりうる入力があった(模した版で確かめた)。直して、渡した先で照らす自己診断を付けた。
\item 印 3 は 全体の 6.1--6.7\% から 1.7--2.2\% に減った。$\mu=10^{-13}\eV$ の軽い BH が「模型の外」になる問題はなくなった。印 5 は 2.6\% から 3.5--3.6\% に増えた(1 世代目の層で 0.9--1.1\%、GW241011 の層で 8.4--9.1\%)。
\item 単体版は pickle でき、メモリは 0.4 GB、警告を出さない。
\end{itemize}
\end{result}

\begin{caveat}
\begin{itemize}
\item 約束は、標本で確かめた統計的な主張である。印 0・1・2 を返す独立な 4{,}700 点で越えた点が 0 なので、点あたりの違反の割合は 95\% の信頼度で $6\times10^{-4}$ 以下と言える(同じ点の年齢どうしは独立でないので、点の数で数える)。
\item 約束は基準の運動論への忠実さである。基準の運動論の率そのものの誤差(解析式の 3--23\%、322 の表の節の 9\%)は残っている。率の表を多倍長で作り直し、値の表を作り直せば、所見 1 の折れと印 5 の増加分はなくなる。未着手。
\item 折れの判定は $\alpha$ の方向だけである。$s$ と $f$ の方向は E18 の曲がりの指標(0.05 dex 以上)のまま。照合で越えるセルがないことは確かめたが、同じ種類の小さな折れがないことを表から直接示してはいない。
\item 平らな所の値は、角の節のうち進みのいちばん遅い節で決まる。角の節の値の幅が 0.02 以下の点は印がつかないので、その幅までのずれは残る。
\item 422 は印 4 の判定に入っていない。(211, 322) の固定点では主要項で育たない(余裕は 2.1 倍)。(322, 544) の固定点では確かめていない。
\item x86 の実機では走らせていない。
\item 固定点で $\alpha>0.30$ の点は印 3 のまま。速く回る重い BH と、$\mu\gtrsim1.5\times10^{-12}\eV$ での $40\Msun$ より重い BH は、大半が約束の外にある(表~\ref{tab:e19cov})。
\item 査読で見ていないもの: 2026 年 7 月以降の文献の数値(\arxiv{2607.01317}、\arxiv{2609.02678} など。原文がこの計算機にない)、E8--E10 の天体物理の入力、勉強ノートの本文、基準の表を作るスクリプトの再実行(出力の表を照合した)。
\end{itemize}
\end{caveat}
